% Source: https://artofproblemsolving.com/community/c3h2183777_stanford_contest_19461965_problem_set
% Stable identity: aops-topic-2183777-q02
\textbf{Method 1.} For every odd $k=2m+1$,
\[x^k-x=x\left(x^{2m}-1\right)=x\left(\left(x^2\right)^m-1\right),\]
and $x^2-1$ divides $\left(x^2\right)^m-1$ (since $t-1$ divides $t^m-1$). So $x(x^2-1)=x^3-x$ divides $x^k-x$. All five exponents $1,9,25,49,81$ are odd, so
\[x+x^9+x^{25}+x^{49}+x^{81}-5x=\left(x^9-x\right)+\left(x^{25}-x\right)+\left(x^{49}-x\right)+\left(x^{81}-x\right)\]
is divisible by $x^3-x$. Since $5x$ has degree less than $3$, the remainder is $\boxed{5x}$.

\textbf{Method 2 (remainder theorem).} Write $P(x)=(x^3-x)Q(x)+R(x)$ with $R(x)=\alpha x^2+\beta x+\gamma$. The divisor vanishes at $0$, $1$ and $-1$, so $R(0)=P(0)=0$, $R(1)=P(1)=5$ and $R(-1)=P(-1)=-5$. Then $\gamma=0$, $\alpha+\beta=5$ and $\alpha-\beta=-5$, so $\alpha=0$, $\beta=5$, and again $R(x)=5x$.
